\section{Discussion}
\label{sec:discussion}

Extending the single-patch model of \citep{RodriguezFalcon2026} to a network of patches changed less than the network literature might lead one to expect, as long as dispersal was free. The amount of parasitoid released, and therefore the release strategy, determined the pest load far more than topology did, and the degree of the outbreak patch, although a nearly deterministic predictor, moved the network AUC by less than one per cent. Transit mortality changed this picture. Once parasitoids died while moving, the efficiency of targeted releases acquired a pronounced maximum at an intermediate dispersal rate, its position scaled with the transit losses summed over the whole network, and two of the structural results obtained without mortality, the equivalence of the star and the complete graph and the benefit of a well-connected outbreak patch, no longer held.

\subsection*{Release strategy, network structure and the outbreak patch}

Uniform releases lowered the network pest load by about 30\% and targeted releases by 2.5--4.5\%, but a uniform programme used eight times as many parasitoids. Per parasitoid released the two strategies were comparable, and the choice between them depended on how fast each species moved. Concentrating releases at the outbreak patch paid off once the parasitoid was mobile enough. With a background infestation (IC-A) this required $\Dt$ above about 0.05 for any $\Dd$ between $10^{-3}$ and 1, and only after a point introduction (IC-B), and for $\Dd < 0.03$, did the threshold follow the dispersal rate of the psyllid. In both cases targeting was clearly worse when psyllids moved and parasitoids did not. The comparison with network epidemiology is instructive here. In epidemic models the degree of the seed node and the targeting of hubs can decide whether a perturbation invades the network \citep{PastorSatorras2001,Colizza2007}, and dispersal network structure can alter metapopulation dynamics strongly even in simple ecological models \citep{holland2008,gilarranz2012}. In our model the corresponding effects were consistent but small. A likely reason is that every patch carries the same local dynamics and that the dispersal rates of Experiments~2 and~3, at most 0.1, couple the patches weakly over the horizon of an outbreak, so topology can only modulate how fast the initial asymmetry is redistributed. Proposition~\ref{prop:hub} points to where the structural information that does matter resides. When the outbreak patch is adjacent to all others, the edges among the remaining patches are dynamically irrelevant. The proposition covers only that case, but across the 34 graphs of Experiment~3 the source degree also summarised how directly the outbreak was coupled to the rest of the network.

That summary is fragile. The negative association between source degree and pest load required a background infestation (IC-A), dispersal rates near $0.03$, and no transit mortality. Under a point introduction (IC-B) a well-connected outbreak patch spread a pest that was absent elsewhere and the association became positive, and with parasitoid transit mortality the sign reversed already at $m_t \approx 0.2$.

\subsection*{Spatial coherence over the horizon of an outbreak}

The master stability function \citep{pecora1998,arenas2008} predicted that the synchronous state is linearly stable on all seven topologies, and the ranking of the transverse exponents broadly agreed with the ranking of the observed synchrony. Their magnitudes, however, imply e-folding times between 60 and about 1050 time units, against a simulated horizon of 100, so over an outbreak a transverse perturbation shrinks only by a factor between 0.19 and 0.91. The star and the complete graph, whose exponents differ by a factor of about seven, had identical synchrony. The linear analysis therefore orders the networks but does not set the level of coherence, which depended on the initial asymmetry and on the way the outbreak patch was coupled to the rest, a transient property rather than an asymptotic one \citep{hastings2018}. The exponents are also hard to estimate. Over a trajectory of 840 time units they were of the same order as their sampling error, and the path even appeared unstable, so only a much longer trajectory resolved their signs. The local dynamics add to this, because the isolated patch oscillates with a period close to 7.8 time units and periodic releases did not entrain it. With patch-level oscillations that are locked neither to the release calendar nor, over an outbreak, to each other, a familiar setting for spatial synchrony in natural populations \citep{liebhold2004}, the asymptotic stability of a synchronous orbit is a weak guide to an outbreak of finite duration.

\subsection*{Transit mortality and an optimal dispersal rate}

\citet{heimpel2011} proposed that biological control agents should disperse neither too little nor too much. Our results support a restricted version of this idea. Without a cost of dispersal there was no intermediate optimum. Parasitoids that left the outbreak patch were not lost, since moving carried no cost beyond natural mortality, and the efficiency of targeted releases saturated at high dispersal rates instead of declining. With transit mortality, a cost documented for many dispersing organisms \citep{bonte2012}, a maximum appeared in a third of the curves, with a median prominence of 73\%, and moved to lower $\Dt$ almost as $1/m_t$.

The maximum is not a minimum of the pest load, and this qualification matters for how the result is used. With transit mortality the pest load under targeted releases tended to rise with $\Dt$. Its minimum was usually at or near the lowest dispersal rate, although on the path and the cycle after a point introduction interior minima reached 15\% below the lower end of the curve. The optimum arose because the no-release reference deteriorated. As $\Dt$ increased, transit losses drove the resident parasitoid population towards extinction and the pest load without releases more than doubled on the complete graph, while repeated releases still sustained parasitoids near the outbreak patch. Each released parasitoid then avoided a large pest load. The optimum identifies the conditions under which releases make the largest difference relative to doing nothing, and those are conditions under which natural control has already failed. An efficiency defined against a counterfactual can therefore be highest in landscapes that are the worst for the parasitoid, and comparisons of release programmes across sites should report the pest load achieved together with the efficiency.

The scaling of the optimum with the mean degree $\langle k\rangle$, rather than with the degree of the outbreak patch, and its unexpected dependence on $r_3$ can both be understood from a linear invasion argument. Near $x_3 = 0$ the crowding term of equation~\eqref{eq:local} vanishes, and if the psyllid density $x_1(t)$ is the same in every patch, the parasitoid vector obeys $\dot{\mathbf{x}}_3 = [r_3 + b_{31}x_1(t)]\,\mathbf{x}_3 - \Dt\,(K - (1-m_t)A)\,\mathbf{x}_3$. The first term is a scalar multiple of the identity and commutes with the second, so a resident parasitoid population at low density grows or declines at the rate
\begin{equation}
  \lambda = r_3 + b_{31}\,\bar{x}_1 - \Dt\,\sigma, \qquad
  \sigma = \lambda_{\min}\bigl(K - (1-m_t)A\bigr),
  \label{eq:invasion}
\end{equation}
where $\bar{x}_1$ is the time average of the psyllid density in the absence of the parasitoid. Because the uniform vector is the null vector of $L$, first-order perturbation theory gives $\sigma = m_t\langle k\rangle + O(m_t^2)$, and the Rayleigh quotient gives $\sigma \le m_t\langle k\rangle$ exactly, with equality on regular graphs. For the eight graphs of Experiment~10, $\sigma/(m_t\langle k\rangle)$ lies between 0.91 and 1 for $m_t \le 0.1$ and between 0.66 and 1 for $m_t = 0.4$. Setting $\lambda = 0$ yields the threshold $\Dt m_t \langle k\rangle \approx b_{31}\bar{x}_1 - |r_3|$ suggested empirically in section~\ref{sec:res-rule}, which gives 0.50 and 0.35 for $r_3 = -0.3$ and $-0.45$ under IC-B, against 0.46 and 0.35 observed, and 0.56 and 0.41 under IC-A, against 0.49 and 0.39. Because $\sigma \le m_t\langle k\rangle$, equation~\eqref{eq:invasion} places the collapse of the resident parasitoid at or above these values, yet the optima lay slightly below them although the resident was already declining there. The likely reason is that $\bar{x}_1$, measured once the parasitoid has gone, overestimates the host density available to the parasitoid while it is still present and suppressing the psyllid, which lowers the effective threshold. Equation~\eqref{eq:invasion} also explains why a larger $|r_3|$ lowers the tolerable transit loss, since the parasitoid has less margin for growth before dispersal costs are added. With edge-specific mortalities the same first-order argument gives $\sigma \approx \langle k\rangle\,\bar{m}$, with $\bar{m}$ the mean of $m_{ij}$ over all edges, which is the mean-degree form that predicted the optimum of Experiment~11 within 10\%. Psyllid transit mortality lowers $\bar{x}_1$, and with it the threshold, which is consistent with the shift of $\Dt^*$ to lower values in Experiment~10b. The argument is a heuristic. It assumes a spatially uniform psyllid density, which IC-B violates early in the outbreak, and it locates the collapse of the resident population rather than the efficiency maximum itself, which also depends on the finite horizon. It predicts, however, that on irregular graphs $\rho_k$ should drift upwards at large $m_t$, where $\sigma$ falls below $m_t\langle k\rangle$, and this can be checked against the stored output. The quantity $\sigma$ plays the same role as the leading eigenvalue that defines the metapopulation capacity of a landscape \citep{hanski2000}, a landscape-level property that combines topology and movement costs into a single persistence threshold.

The failure of every dimensionless group to collapse the efficiencies in Experiment~8 points in the same direction. Most of the variance that dispersal rates could not explain was due to $r_3$, which changes how much pest there is to remove and not only how long a parasitoid lives. A predictive scaling variable would have to include a measure of the pest load available to be suppressed, such as $\mathrm{AUC}_0$ or $b_{31}\bar{x}_1/|r_3|$, alongside the dispersal time scales.

\subsection*{Implications for release programmes against \emph{Diaphorina citri}}

Augmentative biological control is often judged on cost per unit of pest suppressed \citep{collier2004,vanlenteren2012}, and three results bear on that judgement. First, efficiency fell with release size for every period, so each additional parasitoid removed less pest load than the previous one. The most efficient schedule lay at the edge of the grid and achieved reductions below 7\%, so the efficiency optimum is not an operating point, and programmes have to trade suppression for efficiency explicitly. Second, targeting the outbreak patch was competitive per parasitoid only above a threshold of parasitoid mobility, which after a point introduction rose with the mobility of the psyllid. \emph{D.~citri} can move between groves and over long distances \citep{boina2009,lewisrosenblum2015}, whereas the dispersal radius of released \emph{T.~radiata} measured in Brazilian groves was about 18~m \citep{marin2023}. Model rates are in arbitrary units and cannot be matched directly to these measurements, but if released parasitoids stay within about 20~m of the release point while psyllids move between groves, targeted releases into a single focus are likely to lose to area-wide releases, above all after a point introduction, which agrees with the emphasis on coordinated regional management of HLB \citep{bassanezi2020}. Third, the efficiency of releases into a landscape that is costly to cross will look high precisely because resident parasitoids cannot persist there, and the same landscape will carry a high pest load.

Release accuracy and host-dependent searching also matter, as other spatially explicit release studies have stressed \citep{ARCASIM1996,Diatraea2024,SITMetapopulation2026}. \emph{Telenomus podisi} left host-rich plants more slowly, and releases placed on infestation foci raised parasitism \citep{PatchDispersal2026}. Our parasitoids dispersed at a constant rate irrespective of host density. Host-dependent emigration would retain parasitoids in infested patches and reduce transit losses, which would probably move the optimum to higher $\Dt$ for real parasitoids with directed search.

\subsection*{Limitations}

The model is formulated in arbitrary units and its parameters were not estimated from field data, so our conclusions are qualitative. The networks had eight patches and dispersal was linear and density-independent. Larger landscapes, distance-dependent rates and host-dependent movement \citep{comins1992,RefugeConnectivity2022} could restore a larger role for topology. The model is deterministic, so psyllid densities that fell to $10^{-3}$ under intense schedules would correspond to local extinction in a stochastic model, and our results say nothing about eradication. The sensitivity analysis was local and one-at-a-time and explored only a neighbourhood of the reference point \citep{saltelli2019}. Several conclusions depended on the initial condition, and the horizon of 100 time units, chosen to represent an outbreak, is short relative to most transverse time scales and to the growth of the intrinsic oscillation. Finally, the crowding coefficients of the local model can become negative, so boundedness relies on interaction terms, a feature inherited from the single-patch formulation \citep{RodriguezFalcon2026} that limits the range of release schedules over which the model can be trusted.
